% =====================================================================
%  From Stack Smashing to Code Reuse  --  Final Project Report
%  Two-column conference-paper style on the stock `article` class.
%  (IEEEtran is not required; this compiles with a standard TeX Live.)
%
%  Build:   pdflatex report.tex && pdflatex report.tex
%     or:   latexmk -pdf report.tex
% =====================================================================
\documentclass[10pt,twocolumn,letterpaper]{article}

\usepackage[letterpaper,margin=0.75in,columnsep=0.28in]{geometry}
\usepackage{newtxtext,newtxmath}   % Times text + math
\usepackage[T1]{fontenc}
\usepackage[utf8]{inputenc}
\usepackage{microtype}
\usepackage{booktabs}
\usepackage{array}
\usepackage{tabularx}
\usepackage{enumitem}
\usepackage{xcolor}
\usepackage{listings}
\usepackage{caption}
\usepackage[labelsep=period]{caption}
\usepackage{titlesec}
\usepackage{balance}
\usepackage[hidelinks]{hyperref}

% ---- IEEE-ish sectioning: I. SECTION, A. Subsection -------------------
\renewcommand{\thesection}{\Roman{section}}
\renewcommand{\thesubsection}{\Alph{subsection}}
\titleformat{\section}{\centering\normalsize\bfseries\scshape}{\thesection.}{0.5em}{}
\titleformat{\subsection}{\normalsize\itshape}{\thesubsection.}{0.5em}{}
\titlespacing*{\section}{0pt}{10pt}{4pt}
\titlespacing*{\subsection}{0pt}{6pt}{2pt}

% ---- table / caption look --------------------------------------------
\captionsetup{font=small,labelfont=sc}
\captionsetup[table]{position=top}
\setlist[itemize]{leftmargin=1.2em,itemsep=1pt,topsep=2pt,parsep=0pt}

% ---- code listings ----------------------------------------------------
\definecolor{cmt}{RGB}{90,110,90}
\lstset{
  basicstyle=\ttfamily\footnotesize,
  commentstyle=\color{cmt},
  breaklines=true,
  columns=fullflexible,
  keepspaces=true,
  frame=single,
  rulecolor=\color{black!30},
  backgroundcolor=\color{black!3},
  xleftmargin=2pt,xrightmargin=2pt,
  aboveskip=4pt,belowskip=2pt,
}
\lstdefinestyle{py}{language=Python,morecomment=[l]\#}
\lstdefinestyle{c}{language=C}

\newcommand{\code}[1]{\texttt{\small #1}}

% =====================================================================
\begin{document}

% ---- Full-width title band (spans both columns) ----------------------
\twocolumn[{%
  \begin{center}
    {\LARGE\bfseries From Stack Smashing to Code Reuse:\\[2pt]
     A Reproducible Study of the Stack Buffer Overflow\\ and the NX, ASLR, and Canary Defenses\par}
    \vspace{8pt}
    {\large Al Monaem Khan \quad Juan Reynosa \quad David Duru\par}
    \vspace{3pt}
    {\itshape Department of Computer Science, Computer Security (Fall 2026)\par}
    \vspace{2pt}
    {\small\url{https://github.com/al-monaem/computer-security-project}\par}
    \vspace{14pt}
  \end{center}
}]

% ---- Abstract + index terms begin column 1 ---------------------------
\noindent\textbf{\textit{Abstract}---We present a reproducible teaching study
that follows one vulnerability class, the stack buffer overflow, through its
history of attack and defense. We attack a single deliberately vulnerable C
program three ways, changing only its compile-time defenses between runs. In
Stage~1 we inject machine code onto an executable stack and redirect control
into it (x86-64). In Stage~2 we enable the NX bit, show the same exploit
faulting, and recover execution with a return-to-libc call into
\code{system("/bin/sh")}. In Stage~3 we keep NX enabled and build a
Return-Oriented Programming (ROP) chain that assembles an \code{execve} syscall
by hand from four gadgets, then extend it to defeat ASLR with an information
leak and to re-enable execution at runtime through an \code{mprotect} chain
(x86-32). Because only one variable changes between stages, each result isolates
the effect of a single mitigation. The experiments show that no one defense is
enough on its own, which is why modern systems combine several.}

\vspace{4pt}
\noindent\textbf{\textit{Index Terms}---}\textit{stack buffer overflow,
control-flow hijacking, NX/DEP, return-to-libc, return-oriented programming,
ASLR, defense-in-depth.}
\vspace{6pt}

% =====================================================================
\section{Introduction}

\subsection{Problem Statement}
Memory-corruption bugs are among the oldest and most studied software
vulnerabilities, and the stack buffer overflow is the canonical
example~\cite{alephone}. When a program copies attacker-controlled input into a
fixed-size stack buffer without checking its length, the excess bytes overwrite
adjacent stack data, most importantly the saved return address. By controlling
that address, an attacker redirects the program's control flow to code of their
choosing. This work asks a focused question: given one such bug, what does it
take to exploit it, and what does each modern defense actually prevent?

\subsection{Motivation}
Stack overflows have caused serious real-world incidents. The 1988 Morris Worm,
the first large-scale internet worm, spread in part by overflowing a fixed-size
stack buffer in the Unix \code{fingerd} daemon, which read network input with
the unbounded \code{gets()} function~\cite{alephone}; it reached roughly 10\% of
the ${\sim}60{,}000$ hosts then online and led to the first CERT Coordination
Center. The same bug class was still doing damage decades later: the 2017
EternalBlue exploit (CVE-2017-0144) triggered a buffer overflow in Microsoft's
SMBv1 service~\cite{cve}, and the WannaCry ransomware built on it reached more
than 200{,}000 systems across about 150 countries. NX, ASLR, and stack
canaries, the defenses studied here, were introduced in response to attacks
like these~\cite{pax,stackguard}.

\subsection{Contributions}
This paper makes the following contributions:
\begin{itemize}
  \item A single vulnerable program attacked across three configurations, so
    each experiment isolates the effect of exactly one defense (Table~\ref{tab:stages}).
  \item A verified classic injection on an executable stack with a
    runtime-measured offset (Stage~1).
  \item A demonstration that the NX bit converts the same exploit into a crash,
    and that return-to-libc routes around it (Stage~2).
  \item A verified ROP suite (a four-gadget \code{execve} chain, an ASLR
    information-leak bypass, and an \code{mprotect} runtime W\^{}X bypass), plus
    an analysis of why each residual defense would have raised the bar (Stage~3).
\end{itemize}

\subsection{Scope}
The work targets x86 Linux. Stages~1 and~2 use an x86-64 build of the target;
the Stage~3 ROP prototype uses a 32-bit (x86-32) build so that syscall arguments
are passed on the stack, keeping the gadget layout simple to teach. All testing
is confined to an isolated SEED Ubuntu VM~\cite{seed} with harmless payloads.
Out of scope: networked or third-party targets, weaponized payloads, heap
corruption, and non-x86 architectures.

% =====================================================================
\section{Background and Related Work}
The foundational references map directly onto our stages. Aleph One's
``Smashing the Stack for Fun and Profit''~\cite{alephone} documents the classic
injection attack reproduced in Stage~1. Solar Designer's return-to-libc
technique~\cite{solar} is reproduced in Stage~2, and Shacham's generalization to
Return-Oriented Programming~\cite{rop} is the basis of Stage~3. The defenses
correspond to the PaX project's NX and ASLR work~\cite{pax} and to StackGuard
for canaries~\cite{stackguard}. The SEED Labs~\cite{seed} provide the standard
pedagogical environment we build on.

Existing tutorials typically demonstrate one technique against one binary. Our
contribution is methodological: we attack the \emph{same} program across all
stages and change only the build flags between them, so each result is a
controlled before/after that isolates a single mitigation, which most
single-technique writeups cannot show. A few terms recur below: a \emph{gadget}
is a short existing instruction sequence ending in \code{ret}; a ROP
\emph{chain} strings gadget addresses on the stack so each \code{ret} invokes
the next; W\^{}X / NX means a page is writable or executable but not both.

% =====================================================================
\section{Threat Model and Security Requirements}

\subsection{System Overview}
The target is \code{review\_intake.c}, a small command-line ``Product Review
Intake'' tool. Its \code{submit\_review()} function declares a 64-byte
\code{comment} buffer and fills it with an unchecked
\code{read(fd, comment, 256)}. This is one believable defect: a buffer sized at
64 bytes but read up to 256. The same source is compiled under different flags
per stage.

\subsection{Assets and Adversary}
\textbf{Assets.} The integrity of the process's control flow (saved return
addresses and instruction pointer) and, in turn, the confidentiality and
availability of the host. \textbf{Threat actor.} An attacker who can supply
input to the program but holds no legitimate code-execution privilege and, in
the baseline case, cannot read the victim's memory map directly. In the ASLR
scenario the attacker additionally observes program output, which enables the
information leak.

\subsection{Attack Surface and Threats}
\begin{itemize}
  \item Unvalidated input copied into a fixed-size stack buffer (the
    overflow), which is the entry point.
  \item The saved return address as the control-flow hijack point.
  \item Injected shellcode on an executable stack; once NX forbids that,
    existing libc code reused as a single function (ret2libc) or as chained
    gadgets (ROP).
\end{itemize}

\subsection{Assumptions and Requirements}
We assume the attacker controls only program input and that the defender
compiles with standard mitigations; for teaching, specific mitigations are
disabled per stage and ASLR is toggled at the OS level. The security
requirement is an \textbf{integrity} property on control flow: attacker input
must not lead to execution of attacker-chosen behavior. That property underpins
the confidentiality and availability of the host. We examine how NX, ASLR, canaries,
and PIE each contribute, and where each fails alone.

% =====================================================================
\section{Design, Methods, and Implementation}

\subsection{Overall Design}
One vulnerable program is paired with a tool per stage that runs it under a
specific defense configuration. Compiler and linker flags turn the relevant
protection on or off, so the binary is deterministic and each stage isolates a
single variable. The three stages form an attack/defense arc, summarized in
Table~\ref{tab:stages}.

\begin{table*}[t]
  \centering
  \caption{The three stages attack the same program; only the defense changes between them.}
  \label{tab:stages}
  \small
  \begin{tabularx}{\textwidth}{@{}c l X X@{}}
    \toprule
    \textbf{Stage} & \textbf{NX (stack)} & \textbf{Technique} & \textbf{Result}\\
    \midrule
    1 & OFF & Inject shellcode onto the stack, return into it (x86-64) & Shell / hijack confirmed\\
    2 & ON  & Same exploit re-run $\rightarrow$ fails; then return-to-libc into \code{system("/bin/sh")} & SIGSEGV, then shell via reuse\\
    3 & ON  & Full ROP chain (\code{execve} from 4 gadgets); + ASLR leak; + \code{mprotect} (x86-32) & Shell despite NX and ASLR\\
    \bottomrule
  \end{tabularx}
\end{table*}

\subsection{Tools and Technologies}
\begin{itemize}
  \item \textbf{GCC} (with \code{-m32} for the 32-bit target) on an isolated
    SEED Ubuntu 20.04 VM (glibc 2.31).
  \item \textbf{GDB} for stack inspection and step-by-step verification of the
    hijack; \code{readelf}/\code{objdump} for binary inspection.
  \item \textbf{pwntools}~\cite{pwntools} for exploit scripting and
    \textbf{ROPgadget}~\cite{ropgadget} for gadget discovery; 32-bit toolchain
    via \code{gcc-multilib}, \code{libc6-dev-i386}.
\end{itemize}

\subsection{Build Configurations}
Each stage compiles the identical source with flags that leave every protection
on except the ones that would make the lesson impossible to observe
(Table~\ref{tab:builds}). Critically, only one flag
(\code{-z execstack}~$\rightarrow$~\code{-z noexecstack}) separates Stage~1 from
Stage~2.

\begin{table*}[t]
  \centering
  \caption{Per-stage build. Only the NX flag differs between Stage~1 and Stage~2; ASLR is toggled at the OS level (\code{setarch -R} / pwntools).}
  \label{tab:builds}
  \small
  \begin{tabularx}{\textwidth}{@{}c c c c c X@{}}
    \toprule
    \textbf{Stage} & \textbf{Arch} & \textbf{NX (stack)} & \textbf{Canary} & \textbf{PIE} & \textbf{Key flags}\\
    \midrule
    1 & x86-64 & OFF & off & off & \code{-z execstack, -fno-stack-protector, -no-pie}\\
    2 & x86-64 & ON  & off & off & \code{-z noexecstack} (one flag flipped)\\
    3 & x86-32 & ON  & off & off & \code{-m32, -z noexecstack, -fno-stack-protector, -no-pie}\\
    \bottomrule
  \end{tabularx}
\end{table*}

\subsection{The Planted Vulnerability}
The single vulnerability, identical across all stages:
\begin{lstlisting}[style=c]
static void submit_review(void) {
  char comment[64];     /* fixed buffer: the bug */
  ...
  int n = read(STDIN_FILENO, comment, 256);
}
\end{lstlisting}
\noindent Using \code{read()} (not \code{scanf}/\code{fgets}) matters: it copies
raw bytes of any value, including NUL, so arbitrary addresses survive the copy.

\subsection{Stage 1: Classic Injection (x86-64)}
Built with \code{-z execstack}, so the stack is mapped \code{rwxp}. The exploit
avoids guessing the buffer's runtime address: it launches the target under GDB,
breaks inside \code{submit\_review}, and reads the live address of
\code{comment} and the live saved-return-address slot from the very process it
then exploits. The measured offset is 88 bytes (\code{comment} at
\code{[rbp-0x50]}, return address at \code{[rbp+8]}). The payload is
\code{execve} shellcode, NOP padding, then the little-endian address of the
buffer itself:
\begin{lstlisting}[style=py]
payload  = shellcode              # execve("/bin/sh")
payload += b'\x90'*(off-len(sc))  # pad to offset 88
payload += p64(comment_addr)      # saved RA -> buffer
\end{lstlisting}
\noindent GDB confirmed RIP entered the injected bytes.

\subsection{Stage 2: NX Blocks Injection; ret2libc Recovers It (x86-64)}
\fbox{\parbox{0.95\columnwidth}{\footnotesize\textbf{Status (placeholder).} The
Stage~2 implementation and its captured evidence (the Stage~1 payload re-run
under NX, GDB output, and the finalized x86-64 ret2libc run) are being finalized
for the demonstration. The design below is settled; captured results will be
added to Table~\ref{tab:results}.}}

\vspace{4pt}
Stage~2 rebuilds the same source with a single flag flipped,
\code{-z noexecstack} (NX on), and shows two things. \textbf{(a) The defense
works:} the exact Stage~1 payload now faults with \code{SIGSEGV} the instant the
CPU tries to fetch the injected bytes. NX is enforcing W\^{}X here, with no
source change. \textbf{(b) Code reuse defeats NX:} instead of injecting code, the
return address is redirected into an existing libc function, with a fake frame
supplying its argument:
\begin{lstlisting}[style=py]
payload  = b'A'*OFFSET   # pad to saved RA
payload += p32(system)   # ret INTO system()
payload += p32(exit)     # system()'s return addr
payload += p32(binsh)    # arg -> "/bin/sh"
\end{lstlisting}
\noindent Nothing on the stack executes as code, so NX never fires. (Shown in
the 32-bit argument layout; the x86-64 variant loads the argument into RDI via a
\code{pop rdi; ret} gadget.) This is the split the project turns on: NX stops
\textbf{injection} but not \textbf{reuse}, which is what motivates the general
technique in Stage~3.

\subsection{Stage 3: Full ROP, ASLR Bypass, mprotect (x86-32)}
Stage~3 uses a 32-bit build (stack-passed syscall arguments keep the layout
teachable). The offset, 84 bytes, is found with a cyclic De Bruijn pattern. The
main exploit builds the \code{execve("/bin/sh",0,0)} syscall by hand from four
gadgets:
\begin{lstlisting}[style=py]
payload  = b'A'*OFFSET
payload += p32(pop_ecx_edx)+p32(0)+p32(0) # ecx,edx=0
payload += p32(pop_ebx)+p32(binsh)        # ebx=&sh
payload += p32(pop_eax)+p32(0xb)          # eax=11
payload += p32(int80)                     # int 0x80
\end{lstlisting}
\noindent Each \code{ret} pops the next gadget's address; every instruction that
runs lives inside libc, which is legitimately executable, so NX has nothing to
object to. Two further prototypes extend the idea:
\begin{itemize}
  \item \textbf{Defeating ASLR.} With ASLR on, a first stage calls
    \code{puts()} on a GOT entry to leak a real libc address, computes the base
    (\code{libc.address = leaked\_puts - sym['puts']}), and sends the chain
    resolved against it. It succeeds on repeated runs, each with a different
    base.
  \item \textbf{Runtime W\^{}X bypass.} A ROP chain calls \code{mprotect()} to
    mark a page RWX, then jumps into shellcode placed there. This shows NX can
    be undone at runtime given the right callable function.
\end{itemize}

% =====================================================================
\section{Experimental Evaluation and Results}

\subsection{Setup and Methodology}
All tests run on an isolated SEED Ubuntu 20.04 VM (glibc 2.31), x86-64 host with
32-bit multilib. ASLR is disabled per-process where a deterministic address is
needed (\code{aslr=False} / \code{setarch -R}) and explicitly enabled for the
ASLR-bypass test. Each binary's protections are confirmed with
\code{make checksec} before exploitation. We (i) confirm the build, (ii) measure
the offset empirically rather than assuming it equals the buffer size, and (iii)
deliver the payload and verify the outcome: register/RIP state in GDB for a
hijack, a live shell (\code{id}, \code{uname -m}) for success, or the expected
\code{SIGSEGV} for the defense. Results are summarized in
Table~\ref{tab:results}.

\begin{table*}[t]
  \centering
  \caption{Results summary. Stage~1 and all Stage~3 experiments are verified in the VM; Stage~2 is designed and its captured evidence is being finalized.}
  \label{tab:results}
  \small
  \begin{tabularx}{\textwidth}{@{}l l X l@{}}
    \toprule
    \textbf{Experiment} & \textbf{Config} & \textbf{Outcome} & \textbf{Status}\\
    \midrule
    S1: inject shellcode on stack & NX off, x86-64 & Offset 88; GDB confirms RIP in injected code (test: RAX=0x2a); hijack & Verified\\
    S2(a): re-run S1 payload & NX on, x86-64 & Faults with SIGSEGV; NX blocks stack execution & In progress\\
    S2(b): return-to-libc & NX on & One jump into \code{system("/bin/sh")} yields a shell & In progress\\
    S3: 4-gadget \code{execve} ROP & NX on, x86-32 & Chain assembles \code{execve}; shell obtained; ret2libc cousin also works & Verified\\
    S3: ASLR info-leak & NX on, ASLR on & \code{puts()} leak recovers libc base; succeeds across runs & Verified\\
    S3: \code{mprotect} chain & NX on, x86-32 & Page marked RWX, shellcode runs; shell obtained & Verified\\
    \bottomrule
  \end{tabularx}
\end{table*}

\subsection{Security Evaluation}
Table~\ref{tab:defenses} states, for each defense, what it is and how it bears
on these exploits, including the ones we deliberately left off so the lesson
stays visible.

\begin{table*}[t]
  \centering
  \caption{Each mitigation raises the cost of exploitation but is individually bypassable; only fixing the unbounded write removes the vulnerability.}
  \label{tab:defenses}
  \small
  \begin{tabularx}{\textwidth}{@{}l X X@{}}
    \toprule
    \textbf{Defense} & \textbf{What it is} & \textbf{Effect on these exploits}\\
    \midrule
    NX / DEP & Stack pages non-executable (W\^{}X). & Stops Stage~1 injection (S2a). Does \emph{not} stop reuse; ret2libc and ROP run only existing code.\\
    Stack canary & Random guard before the saved RA, checked on return. & A linear overflow smashes it; the check aborts before \code{ret}, so the RA is never controlled. Would stop all stages as built.\\
    ASLR & Randomizes libc load address each run. & Breaks hard-coded addresses. Bypassed in S3 only via an information leak.\\
    PIE & Randomizes the program image too. & Moves in-binary gadgets, raising the bar for ROP further.\\
    Safe APIs & \code{read(buf, sizeof buf)}; bounds checks. & The overflow never happens; this removes the bug instead of only raising its cost.\\
    \bottomrule
  \end{tabularx}
\end{table*}

\subsection{Analysis}
Flipping one flag turns a working exploit into a crash, which isolates NX's
contribution. Return-to-libc and ROP then recover execution without placing any
code where NX would object, so on its own NX is necessary but not sufficient.
The ASLR experiment adds the next layer: fixed addresses stop being reliable,
and that layer also falls to an information leak, which is why real systems add
canaries and PIE as well. We measured the offsets (88 bytes on x86-64, 84 on
x86-32) rather than assuming them, because the compiler reserves frame space
beyond the 64-byte buffer. None of these mitigations stopped exploitation on its
own, which is why production systems combine them.

% =====================================================================
\section{Discussion, Limitations, and Conclusion}

\subsection{Limitations}
Several mitigations are deliberately disabled per stage for pedagogy; this is not
a hardened production build, and demo addresses are deterministic only because
ASLR is toggled for teaching. The Stage~3 ROP prototype is 32-bit for layout
simplicity rather than matching the x86-64 Stage~1/2 target, and the Stage~2
captured evidence is still being assembled. We do not treat heap corruption,
canaries in practice, jump-oriented programming, or control-flow integrity in
depth.

\subsection{Future Work}
\begin{itemize}
  \item Finalize the Stage~2 evidence and re-target the ROP prototype to the
    x86-64 binary.
  \item Add canary-enabled and PIE-enabled builds to show the overflow caught
    and gadget addresses moving.
  \item Explore jump-oriented programming and a basic CFI demonstration as the
    next defense layer.
\end{itemize}

\subsection{Conclusion}
Attacking one program three ways, changing only its defenses, puts the whole arc
in one place. The classic overflow works when the stack is executable; NX turns
the same attack into a crash; code reuse routes around NX; and ASLR, the next
layer, falls to an information leak. No single mitigation stops exploitation by
itself, so real systems combine them, and the durable fix is to remove the
unbounded write that starts the chain.

% =====================================================================
\begin{thebibliography}{9}
\small
\bibitem{alephone} Aleph One, ``Smashing the Stack for Fun and Profit,''
  \emph{Phrack}, vol.~7, no.~49, 1996.
\bibitem{solar} Solar Designer, ``Return-to-libc Attack,'' Bugtraq mailing
  list, 1997.
\bibitem{rop} H.~Shacham, ``The Geometry of Innocent Flesh on the Bone:
  Return-into-libc without Function Calls (on the x86),'' in \emph{Proc. ACM
  CCS}, 2007.
\bibitem{stackguard} C.~Cowan et al., ``StackGuard: Automatic Adaptive
  Detection and Prevention of Buffer-Overflow Attacks,'' in \emph{Proc. USENIX
  Security}, 1998.
\bibitem{pax} The PaX Team, ``PaX: NOEXEC and ASLR Documentation,''
  \url{https://pax.grsecurity.net/}.
\bibitem{seed} W.~Du, ``SEED Labs: Buffer Overflow and Return-to-libc Attack
  Labs,'' \url{https://seedsecuritylabs.org/}.
\bibitem{cve} MITRE, ``CVE-2017-0144 (EternalBlue / SMBv1 buffer overflow).''
\bibitem{pwntools} Gallopsled, ``pwntools: CTF framework and exploit
  development library,'' \url{https://github.com/Gallopsled/pwntools}.
\bibitem{ropgadget} J.~Salwan, ``ROPgadget: gadget finder,''
  \url{https://github.com/JonathanSalwan/ROPgadget}.
\end{thebibliography}

\balance
\section*{Appendix: Team Contributions}
\begin{itemize}
  \item \textbf{Juan Reynosa}, Stage~1 (classic injection, x86-64): build
    configuration, payload construction, and GDB hijack verification.
  \item \textbf{David Duru}, Stage~2 (NX protection + return-to-libc):
    enabling NX on the x86-64 build, showing the Stage~1 payload fail under NX,
    and the ret2libc workaround.
  \item \textbf{Al Monaem Khan}, Stage~3 (ROP, x86-32): the vulnerable
    application, the four-gadget \code{execve} chain, the ASLR information-leak
    bypass, and the \code{mprotect} chain.
\end{itemize}

\end{document}
